\section{Dataset}
\label{sec:data}

This study rests on a pooled twenty-subject synchronized video--EEG cohort. To our knowledge it is the largest paired rodent video--EEG
corpus studied for learned seizure detection, and its scale, severity labels, and
severely long-tailed burden are what make the head-to-head comparisons in
Sec.~\ref{sec:application} meaningful.

\noindent\textbf{Acquisition.}
Synchronized video--EEG was recorded from twenty rats under chemically induced
seizure protocols. Each animal was filmed in its home cage by a fixed overhead
camera while implanted electrodes streamed synchronized EEG
(Fig.~\ref{fig:setup}). A domain expert annotated every seizure and graded its
behavioral severity on the Racine scale~\cite{racine1972}. Because camera
geometry differs per cage, each clip is cropped per cage before subjects are
pooled.

\noindent\textbf{Clip extraction.}
Around each expert annotation we extract a seizure clip and sample a matched
non-seizure clip from a seizure-free interval of the \emph{same} session,
yielding a pooled corpus of 24{,}131 clips over 601 recording sessions
(12{,}197 non-seizure, 11{,}934 seizure); the same corpus expressed in recorded
time is given in App.~\ref{sec:dataextra:cohort}, Table~\ref{tab:corpusscale}. Sampling the negative from the same
session fixes camera, cage, and lighting across the pair, so a detector must
separate the two from behavior alone (Fig.~\ref{fig:paired} shows one
synchronized pair; the binary task is visualized in App.~\ref{sec:dataextra:stages},
Fig.~\ref{fig:detexample}).

% fig:paired now lives in sec/1_intro.tex so it prints on p.2 alongside its first
% mention in the contributions list.

\noindent\textbf{Severity stages.}
Each clip carries a Racine severity label. The behavioral signal grows more overt
with stage, from mild movement at Stage\,2 to pronounced rearing and convulsive
motion at Stage\,4--5. This is the signal a video model reads directly and a
single EEG channel cannot see, and it foreshadows the grading gap in
Sec.~\ref{sec:application}. Per-stage \emph{filmstrips} illustrating this
progression are provided in the supplement (App.~\ref{sec:dataextra:stages},
Fig.~\ref{fig:stages}).

\noindent\textbf{Long-tailed burden.}
Seizure burden is severely uneven across animals. The four busiest supply nearly
half of all clips, subject totals span nearly three orders of magnitude (11--3{,}091
clips), and the most severe stage (S5, 193 clips in the entire cohort) stays rare
throughout; the full per-subject composition is tabulated and visualized in the
supplement (App.~\ref{sec:dataextra:cohort}, Table~\ref{tab:cohort} and
Fig.~\ref{fig:cohortdist}). This long tail is the central statistical difficulty
of the cohort. A single pooled score is dominated by the few high-burden animals,
so we report per-subject metrics alongside pooled ones and use inverse-frequency
class weighting throughout. It is also why fine-grained severity grading is so hard.

\noindent\textbf{Availability and de-identification.}
Raw recordings cannot be released, as they carry facility- and protocol-level
identifiers. We instead release a fully de-identified cohort description: the
per-subject statistics (App.~\ref{sec:dataextra:cohort}, Table~\ref{tab:cohort}), the task
definitions, and the splitting protocol. All figures in this paper use
de-identified frames with cage-side identifiers cropped or redacted. This makes the
cohort's composition and our protocol reproducible without exposing restricted
data. The public RodEpil benchmark~\cite{rodepil2025}, on which we externally
validate (Sec.~\ref{sec:application}), offers a fully open complement.
